\section{Introduction}
Economic counterfactuals rarely require only one solution of a dynamic problem. A researcher changes an operating benefit, a cost of adjustment, or an outside option and asks whether the policy response is economically significant. Repeating a numerical solve at each parameter value can be costly. Reusing an approximate policy without knowing its value loss can instead confound an economic reversal with approximation error. This problem is especially acute when controls affect both financial resources and the agent's preferences, and when their joint state determines the termination of an operating contract.

This paper develops Neural Bellman Operators (NBO) as a framework for policy computation and contract counterfactuals. The computational objects are separated: a critic approximates continuation value; an actor proposes a feasible control; an improvement calculation checks that proposal; and an independent evaluation measures the induced policy. Neural approximation is useful only insofar as its errors and computational costs can be related to the economic decision being studied. We therefore make the policy-value guarantee, rather than a neural training objective, the central object of the analysis.

The main application is portfolio choice with adjustable risk aversion and first-exit liquidation. It retains consumption, risky investment, costly preference adjustment, two stochastic shocks, and their covariance. The liquidation payoff is an economic primitive, not a numerical transfer that returns wealth to the interior. We show that a constant operating benefit $m$ and a constant liquidation increment $\ell$ influence choices through the single contrast $d=m-\rho\ell$. The separate levels are observationally equivalent along an annuity-normalization direction. Changing $m$ alone changes the contract, whereas shifting both components coherently cannot reverse the portfolio. This distinction separates an operating-duration incentive from an unconditional preference-hedging interpretation.

We then construct a certified policy library over the contract contrast. Independent evaluation supplies each stored policy's baseline benefit, discounted duration, and adjustment effort. Feasible policy values form a lower envelope; certified anchor values form an upper envelope. Statewise switching among stored policies dominates their lower envelope. A finite breakpoint calculation controls the value loss at every state and date for every parameter in an interval, not only at sampled parameter values. A second calculation bounds the value of all opposite-sign first actions. It certifies economic decision regions even when a global welfare tolerance alone would be too coarse to identify a portfolio sign.

The executed experiment uses the positive finite approximation of the stopped economy and includes both the reference and diagnostic action meshes together with the frozen neural proposals. It repairs the substantial feasible deviations of the initial neural fits by using evaluated policy continuations. The resulting library attains a uniform $10^{-3}$ welfare target over the entire displayed contract interval. The initial portfolio is provably negative over one subinterval and positive over another, with an explicitly reported unresolved band between them. Joint changes in time, state, and action resolution retain the endpoint reversal, while their variation near the switch remains visible. These are quantitative conclusions about specified finite economies; they are not inferred from a small differential residual or presented as a Brownian first-passage error estimate.

A coupled stochastic resource-allocation problem provides a separate test of representation and computational reuse. Its shared investment budget binds, shocks are correlated, and the state cost is nonseparable. An independent nonanticipative scenario-tree problem bounds each tested policy's lifetime cost loss. Same-checkpoint deployment comparisons isolate a neural actor's role as an initialization device. A full convex-quadratic continuation fit tests whether a simpler representation attains the same accuracy. Fresh-state query workloads charge for critic training, online improvement, and complete policy evaluation, and compare them with a solver exploiting the model's convex structure. This identifies the workload at which reuse is useful without attributing the entire advantage to the actor.

The general error analysis and the broader economic extensions remain part of the argument. A finite-horizon theorem distinguishes evaluation, improvement, boundary, and transition-operator errors. Exact policy improvement supplies a conditional benchmark, while counterexamples explain why an attracting actor stationary point or a small smooth residual is insufficient. Merton and Epstein--Zin problems test actual evaluation and control updates. Sophisticated temporal selves require a different improvement criterion from ordinary continuation evaluation. A persistent-capacity game tests full unilateral dynamic deviations at every date and state. These exercises serve different purposes and are not combined into a single success count.

\subsection{Related Literature and the Contribution}
Policy improvement and approximate dynamic programming precede neural approximation. \citet{jacka2015} study continuous-time policy improvement under explicit control hypotheses. Our exact-operator proposition records the boundary and limiting requirements needed for that benchmark. The finite-horizon error account is backward propagation through monotone operators; the empirical contribution is to make the relevant error information operational for an economic counterfactual.

The nearest neural approaches also combine value approximation and control improvement. \citet{kim2025} use neural fixed-policy PDE evaluation within policy iteration; \citet{kim2026} develop interior error bounds when boundary information is unresolved. \citet{lee2024} combine policy iteration with deep operator learning, including transfer across terminal functions. \citet{cai2025} use martingale conditions for value and control representations. Direct neural approximation of stochastic controls appears in \citet{han2016}. Neural evaluation, separate parameter blocks, and avoidance of a closed-form argmax are consequently not claims of priority here.

Our policy library is closely related to successor features and generalized policy improvement \citep{barreto2017}, policy switching \citep{chang2021}, and policy caches \citep{nemecek2021}. In particular, the convex anchor upper bound is already available in the policy-cache literature. We use it, rather than present it as a new discovery. The contribution developed here is an economic and computational specialization with three jointly executed requirements: an exact reduction of stopped-contract primitives; a certified, adaptively constructed continuum of counterfactuals using independently evaluated policies; and action-separation bounds that distinguish a portfolio reversal from numerical uncertainty. The application also reports actual repair costs and complete-query workloads, including the simpler alternatives that make the interpretation of neural approximation testable.

The broader computational literature includes projection and recursive methods \citep{judd1998,stokey1989}, adaptive sparse grids \citep{brumm2017}, and neural differential-equation solvers \citep{raissi2019,han2018}. The appropriate comparator exploits the problem's structure. We use Riccati recursion for an LQR benchmark and convex optimization for the constrained resource problem. The convex neural critic is an application of input-convex approximation \citep{amos2017}, not a new architecture. Its role is to represent a continuation function more flexibly than a convex quadratic while retaining a checkable one-period improvement gap.

The following sections define the operators, establish the approximation account, specify the preference economy, develop contract reuse and action certification, and report the computations. The final substantive section treats recursive utility, temporal selves, and capacity competition. The standalone supplement retains the operator proofs, negative controls, equilibrium conditions, and replication details and adds the contract and all-date verification calculations.
